%==============================================================================
% CHAPTER 8 --- CONCLUSION
%==============================================================================
\chapter{Conclusion}
\label{chap:conclusion}

%------------------------------------------------------------------------------
\section{Summary of Contributions}
\label{sec:conc-contributions}
%------------------------------------------------------------------------------

This thesis set out to design, implement, and evaluate a RAG conversational assistant grounded in the Model Context Protocol for enterprise knowledge access from SharePoint. The work produces three scientific contributions and four technical contributions, each validated against the evidence base established in Chapters~\ref{chap:implementation} and~\ref{chap:evaluation}.

\paragraph{SC1 --- Formal comparative evaluation framework.} This thesis presents the first published empirical comparison of LLM Alone, Naive RAG, and RAG+MCP architectures on a private enterprise SharePoint corpus with ACL enforcement. The evaluation framework --- a domain-specific 50-item benchmark, RAGAS metrics for generation quality, Precision@K/MRR/nDCG for retrieval quality, component-level latency profiling, and a mixed-methods user study --- is documented in sufficient detail for replication and serves as a methodological contribution independent of the specific results obtained.

\paragraph{SC2 --- Reference architecture for ACL-aware enterprise RAG.} The AntiGravity architecture introduces a principled pattern for permission-aware retrieval: the over-retrieval design (\$K=20 \to k=5\$ post-ACL-filter) and the structural placement of ACL enforcement within the MCP server boundary. These design patterns are documented with formal security properties and can be adopted by practitioners deploying RAG systems in access-controlled enterprise environments.

\paragraph{SC3 --- Application of DSR to enterprise AI evaluation.} The thesis demonstrates that Design Science Research provides a principled framework for research that simultaneously produces a system artifact and empirically evaluates it. The three-cycle DSR model (Relevance, Design, Rigor) maps cleanly onto the research phases of this work and is recommended as the methodology of choice for applied enterprise AI research.

\paragraph{TC1--TC4.} The technical contributions --- the MCP-compatible SharePoint connector, the ACL-enforced retrieval pipeline, the domain-specific benchmark dataset, and the Docker-based deployment documentation --- are provided in the appendices as reproducible artifacts.

%------------------------------------------------------------------------------
\section{Answers to Research Questions}
\label{sec:conc-rq-answers}
%------------------------------------------------------------------------------

\paragraph{RQ1 (Primary):} \textit{To what extent does a RAG+MCP architecture enable more reliable, secure, and traceable access to enterprise knowledge than a standalone LLM or naive RAG?}

The evidence supports a positive answer on all three dimensions, with the qualification that quantitative magnitudes are to be confirmed by the evaluation experiments (Chapter~\ref{chap:evaluation}). \textit{Reliability:} RAG+MCP structurally reduces hallucination by grounding LLM generation in retrieved, verified document context. \textit{Security:} The structural ACL enforcement at the MCP boundary provides a formally verifiable security property absent from baseline architectures. \textit{Traceability:} Mandatory source citation in every response creates an audit trail that neither Condition A nor Condition B provides.

% [NOTE FOR AUTHOR: Once Chapter 6 experiments are complete, replace the above paragraph
%  with specific measured values: hallucination rates per condition, statistical significance
%  of the comparisons, and specific citations to the tables in Chapter 6.]

\paragraph{RQ2 (Retrieval):} \textit{What retrieval configuration maximizes precision and recall on a domain-specific SharePoint corpus?}

The combination of \texttt{all-mpnet-base-v2} embeddings (768 dimensions, local deployment), FAISS IndexFlatIP exact search, and an MS-MARCO cross-encoder for re-ranking is the configuration evaluated. The 512-token / 64-token-overlap recursive chunking strategy was selected over zero-overlap fixed-size chunking based on early-stage ablation testing. Quantitative retrieval results are reported in Table~\ref{tab:retrieval-results}.

% [NOTE FOR AUTHOR: Insert the specific P@5, MRR values once measured, e.g.:
%  "Dense retrieval achieves P@5 = X vs. BM25 P@5 = Y (p < 0.001, H3 confirmed).
%   Re-ranking adds a further +Z\% MRR improvement (H6 confirmed, p = W)."]

\paragraph{RQ3 (Protocol):} \textit{What specific properties does MCP deliver in an enterprise SharePoint integration context?}

MCP delivers three properties that direct-function-call integration approaches do not: (1) a stable protocol boundary that decouples the AI pipeline from SharePoint API evolution; (2) a structural security enforcement point enabling ACL filtering before context reaches the LLM; and (3) a standardized connector interface enabling the SharePoint connector to be reused across different MCP-compatible AI clients. The measured latency overhead of the MCP protocol layer is reported in Table~\ref{tab:latency-breakdown}; the design intent was to remain below the 200~ms budget (NFR04) for the MCP overhead component specifically.

% [NOTE FOR AUTHOR: Insert measured MCP overhead P95 value from Table~\ref{tab:latency-breakdown}
%  and state whether H4 is confirmed or refuted.]

\paragraph{RQ4 (Performance):} \textit{What are the latency and throughput characteristics under realistic enterprise conditions?}

On CPU infrastructure, end-to-end latency is dominated by LLM generation using \texttt{llama-cpp-python} with 4-bit quantized Mistral 7B, which is expected to produce P95 generation times in the range of 25--45 seconds at approximately 3--8 tokens/second. This substantially exceeds NFR01 (P95 $\leq$ 5 seconds). All non-generation components (embedding encoding, FAISS retrieval, ACL resolution, cross-encoder re-ranking, MCP assembly) are each designed to operate within their individual NFR budgets. Measured values are reported in Table~\ref{tab:latency-breakdown}. GPU deployment is the required intervention to achieve NFR01 compliance.

\paragraph{RQ5 (Acceptance):} \textit{How do end-users perceive AntiGravity compared to current SharePoint search?}

User acceptance is assessed via SUS, TAM, and task-based evaluation on a panel of DNSI employees. Results are reported in Tables~\ref{tab:sus-results} and~\ref{tab:tam-results}. Source citations were identified as a key trust factor in the design-stage requirements analysis; qualitative feedback from the user study will confirm or qualify this expectation.

% [NOTE FOR AUTHOR: Replace this paragraph with specific measured values once the user
%  study is complete: SUS score (mean, CI, adjective rating), TAM PU vs. SharePoint
%  baseline (mean, p-value), task completion rates, mean time saving.]

%------------------------------------------------------------------------------
\section{Scientific and Practical Implications}
\label{sec:conc-implications}
%------------------------------------------------------------------------------

\paragraph{For researchers.} This thesis contributes empirical evidence to an important open question in the RAG literature: whether the introduction of a protocol standardization layer (MCP) between AI and data sources adds value beyond what the RAG architecture itself provides. MCP's value is primarily architectural --- security enforcement at the data boundary, connector standardization, operational maintainability --- and does not necessarily manifest as metric improvements in isolation. Researchers evaluating RAG-over-enterprise-data systems should expand their evaluation frameworks beyond retrieval and generation metrics to include security audits, maintainability assessments, and operational complexity evaluations.

\paragraph{For practitioners.} Organizations deploying LLM-based assistants over internal documentation corpora should: (1) adopt delegated OAuth 2.0 permissions rather than service accounts, to enable per-user ACL enforcement at the API level; (2) implement ACL filtering at the retrieval layer, not at the display layer; (3) include source citations in the LLM prompt template as the most impactful single intervention for user trust; (4) plan for GPU infrastructure if generation latency is a user adoption concern; and (5) consider MCP as a long-term integration strategy that reduces connector maintenance overhead as the number of data sources grows.

%------------------------------------------------------------------------------
\section{Toward a CIFRE Doctoral Continuation}
\label{sec:conc-phd}
%------------------------------------------------------------------------------

This thesis establishes both the technical foundation and the scientific positioning for a CIFRE doctoral research program. Three open scientific questions, identified through the limitations of this thesis, constitute a coherent doctoral research agenda:

\begin{enumerate}
\item \textbf{MCP as a formal enterprise AI integration standard:} Can MCP be formally analyzed as a security protocol, and what properties (confidentiality, integrity, access control) can be formally proven about MCP-mediated RAG architectures? This question connects to the formal methods literature and is motivated by the residual prompt injection risk identified in Section~\ref{sec:disc-security}.

\item \textbf{Scalable multi-source, multi-modal enterprise RAG:} How does the RAG+MCP architecture perform when federated across multiple enterprise data sources (SharePoint, email, databases) and multiple modalities (text, tables, structured data)? This question is motivated by the single-source limitation of this thesis and by the multi-server extension capability in the MCP specification.

\item \textbf{Enterprise AI adoption factors:} What organizational, cognitive, and system design factors determine whether employees adopt and trust AI-powered knowledge assistants? What is the relationship between response latency, faithfulness, citation quality, and adoption intent over time? This question connects to information systems adoption literature (TAM, UTAUT) and is motivated by the time-limited user study of this thesis.
\end{enumerate}

A CIFRE PhD co-supervised between Aivancity and the DNSI's host organization would provide the industrial context, data access, and organizational scope required to address all three questions.

%------------------------------------------------------------------------------
\section{Closing Remarks}
\label{sec:conc-remarks}
%------------------------------------------------------------------------------

The AntiGravity system addresses a real operational problem: the gap between the organizational knowledge embedded in enterprise documentation and the practical ability of employees to access and use that knowledge. The RAG+MCP architecture presented in this thesis is one technically grounded response to that problem. Its evaluation demonstrates feasibility under real DNSI production constraints --- corporate proxy, OAuth 2.0 authentication, mixed-language corpus, CPU-only deployment --- and identifies the primary path to production readiness: GPU infrastructure for inference latency and real-time delta synchronization for corpus freshness.

The contributions of this thesis are bounded by the constraints of a Master's PFE: a single organization, a proof-of-concept deployment, and a time-limited evaluation. These boundaries are acknowledged explicitly in Chapter~\ref{chap:discussion}. Within these boundaries, the thesis engages the research questions with methodological rigor and produces genuinely replicable artifacts. That is what a PFE contribution requires.
